\chapter{Mastectomy}

\section*{Introduction \& Clinical Indications}
A mastectomy is a major surgical procedure involving the complete surgical excision of all glandular breast tissue, typically including the nipple-areolar complex and a variable amount of overlying skin, from one or both breasts. This intervention is primarily performed for the definitive treatment of breast cancer, particularly when the disease is extensive, multifocal, or when breast-conserving surgery is not oncologically feasible or desired by the patient. Beyond its therapeutic applications, mastectomy also serves a crucial role in the prophylactic setting, significantly reducing the risk of developing breast cancer in individuals identified as having a very high genetic or familial predisposition. The primary aim of the procedure is to eradicate existing cancerous cells or prevent their formation, thereby improving long-term survival and quality of life for affected patients. Modern techniques often incorporate immediate or delayed breast reconstruction to restore body image and improve psychosocial outcomes.

\begin{itemize}
  \item Total Mastectomy
  \item Simple Mastectomy
  \item Modified Radical Mastectomy (MRM)
  \item Skin-Sparing Mastectomy (SSM)
  \item Nipple-Sparing Mastectomy (NSM)
  \item Prophylactic Mastectomy
  \item Risk-Reducing Mastectomy
\end{itemize}

The fundamental principle of mastectomy is the complete surgical excision of all glandular breast tissue to remove existing malignancy or to prevent its development in high-risk individuals. This involves meticulously dissecting the breast tissue from the pectoralis major fascia, extending superiorly to the clavicle, inferiorly to the inframammary fold, medially to the sternum, and laterally to the latissimus dorsi muscle. The primary technical goal is to achieve clear surgical margins, ensuring no residual cancerous cells remain in the operative bed, thereby minimizing the risk of local recurrence.

In cases of invasive breast cancer, the procedure's rationale extends to accurate locoregional staging, which often includes an assessment and potential removal of axillary lymph nodes. This is typically performed either through sentinel lymph node biopsy (SLNB) or, if indicated, a formal axillary lymph node dissection (ALND), to determine nodal involvement and guide adjuvant systemic therapy. The extent of skin removal varies based on the type of mastectomy, ranging from a wide elliptical excision including the nipple-areolar complex (NAC) in a simple mastectomy, to preserving most of the skin envelope or even the NAC in skin-sparing or nipple-sparing approaches, respectively. These modern approaches are driven by the rationale to facilitate immediate reconstructive options, improving cosmetic outcomes and patient quality of life without compromising oncologic safety. The choice of technique is tailored to tumor characteristics, patient anatomy, and reconstructive goals, always prioritizing oncologic efficacy.

Early attempts at breast cancer surgery date back to antiquity, with rudimentary excisions described by ancient Egyptians and Greeks. However, these procedures were often associated with high mortality and poor outcomes. It was not until the late 19th and early 20th centuries that systematic and more effective surgical approaches emerged. William Halsted, an American surgeon, revolutionized breast cancer treatment in 1894 with his radical mastectomy, which involved the en bloc removal of the entire breast, pectoralis major and minor muscles, and all axillary lymph nodes. This extensive procedure, though highly morbid due to significant tissue removal and functional impairment, dramatically improved local control rates and remained the standard of care for nearly 70 years.

The mid-20th century saw a paradigm shift, driven by a growing understanding of breast cancer biology and the limitations of radical surgery. Bernard Fisher's groundbreaking randomized clinical trials, notably NSABP B-04 and B-06 in the 1970s and 1980s, demonstrated that less extensive surgeries, such as modified radical mastectomy (preserving the pectoral muscles) and even breast-conserving surgery (lumpectomy) followed by radiation therapy, achieved comparable survival outcomes for many patients with early-stage disease. This evidence-based evolution led to the widespread adoption of modified radical mastectomy as the new standard, and subsequently, to the increasing use of breast conservation, emphasizing both oncologic efficacy and quality of life. Further advancements include skin-sparing and nipple-sparing techniques, and the routine use of sentinel lymph node biopsy, which have further minimized surgical morbidity while maintaining oncologic safety.

Mastectomy is indicated for a variety of clinical scenarios related to breast cancer treatment and prevention.

\begin{itemize}
  \item To treat invasive breast cancer that is multifocal or multicentric, involving multiple distinct tumor foci within the breast, making breast-conserving surgery impractical.
  \item To manage large breast cancers where the tumor size is disproportionate to the breast volume, making breast-conserving surgery cosmetically unacceptable or oncologically unsafe.
  \item To address recurrent breast cancer following prior lumpectomy and radiation therapy, where further breast conservation is not an option due to previous radiation exposure.
  \item To treat inflammatory breast cancer, a highly aggressive form of the disease that often requires mastectomy as part of a multimodal treatment plan after neoadjuvant chemotherapy.
  \item To manage locally advanced breast cancer after neoadjuvant chemotherapy, if residual disease is extensive or breast conservation is not achievable with clear margins.
  \item To provide prophylactic risk reduction for individuals with a very high lifetime risk of breast cancer, such as those with pathogenic BRCA1/2 gene mutations or a strong family history of breast cancer.
  \item To treat extensive ductal carcinoma in situ (DCIS) that is high-grade, involves a large area, or results in persistently positive margins after attempted breast-conserving surgery.
  \item To manage Paget's disease of the nipple when associated with an underlying invasive or extensive in situ carcinoma within the breast.
  \item Patient preference for mastectomy over breast-conserving surgery and radiation, after thorough discussion of all options.
\end{itemize}

Mastectomy serves as a primary surgical treatment for breast cancer in a significant proportion of patients, particularly when breast-conserving surgery (lumpectomy) is not oncologically appropriate or is contraindicated. This includes cases of large tumors, multifocal or multicentric disease, inflammatory breast cancer, or when patients prefer mastectomy to avoid radiation therapy or for personal reasons after comprehensive counseling. Furthermore, it is the primary and most effective risk-reducing option for individuals at extremely high genetic risk of developing breast cancer, such as those with pathogenic mutations in BRCA1, BRCA2, or other high-penetrance genes like TP53 (Li-Fraumeni syndrome) or PTEN (Cowden syndrome).

In some scenarios, mastectomy may be considered a secondary or salvage procedure. For instance, it is performed as a salvage operation for local recurrence after prior breast-conserving surgery and radiation, where further breast conservation is not possible. It may also be performed following neoadjuvant therapy if residual disease precludes breast conservation or if the patient's response to therapy indicates a need for more extensive local control. Its role is determined by a multidisciplinary team, including surgical oncologists, medical oncologists, radiation oncologists, and plastic surgeons, based on tumor biology, patient factors, and shared decision-making, always aiming for the best oncologic and quality of life outcomes.

\section*{Relevant Surgical Anatomy}
The breast is a modified apocrine gland situated on the anterior chest wall, extending from the second to the sixth rib vertically and from the sternal edge to the mid-axillary line horizontally. It is composed of glandular tissue (15-20 lobes, each with multiple lobules and ducts), fibrous connective tissue (Cooper's ligaments, which provide support), and adipose tissue, all encased within the superficial fascia. Deep to the breast lies the pectoralis major muscle, separated by the retromammary space, which facilitates surgical dissection. Other important muscles in the region include the pectoralis minor, serratus anterior, and the latissimus dorsi, which form the boundaries of the axilla.

The arterial supply to the breast is rich, primarily derived from the internal mammary (internal thoracic) artery (medial perforators), the lateral thoracic artery (from the axillary artery), and the thoracoacromial artery. Intercostal perforators also contribute. Venous drainage largely parallels the arterial supply, emptying into the internal mammary, axillary, and intercostal veins. The lymphatic drainage is of paramount surgical importance, with over 75\% draining to the axillary lymph nodes, which are surgically categorized into three levels relative to the pectoralis minor muscle: Level I (lateral to pectoralis minor), Level II (deep to pectoralis minor), and Level III (medial to pectoralis minor). Other lymphatic pathways include the internal mammary nodes and, less commonly, supraclavicular nodes. Key nerves at risk during axillary dissection include the long thoracic nerve (supplying the serratus anterior, injury leads to winged scapula), the thoracodorsal nerve (supplying the latissimus dorsi), and the intercostobrachial nerves (sensory to the axilla and medial upper arm, often sacrificed, leading to numbness). Surgical landmarks include the inframammary fold, the sternal notch, and the axillary tail of Spence, which is an extension of breast tissue into the axilla.

\section*{Preoperative Workup \& Preparation}
Comprehensive preoperative preparation is essential to optimize patient safety, ensure accurate staging, and facilitate successful surgical and reconstructive outcomes. This involves a thorough medical evaluation and specific logistical preparations.

\begin{itemize}
  \item **Medical Evaluation:** A thorough history and physical examination are conducted, along with routine preoperative laboratory tests including a complete blood count, metabolic panel, coagulation studies, and urinalysis. Electrocardiogram (ECG) and chest X-ray may be performed based on age and comorbidities.
  \item **Cardiac and Anesthesia Clearance:** Patients undergo cardiac risk assessment and anesthesia evaluation to ensure fitness for general anesthesia, especially for those with significant comorbidities.
  \item **Imaging and Biopsy Review:** All relevant imaging (mammography, ultrasound, MRI) and pathology reports from prior biopsies are meticulously reviewed to confirm diagnosis, tumor extent, and guide precise surgical planning.
  \item **Genetic Counseling and Testing:** For high-risk patients, genetic counseling and testing for hereditary breast cancer syndromes (e.g., BRCA1/2 mutations) are often performed to inform surgical decisions, including the potential for contralateral prophylactic mastectomy.
  \item **Reconstruction Consultation:** Patients are extensively counseled on all available breast reconstruction options (implant-based, autologous tissue), including the timing (immediate or delayed), and a plastic surgeon consultation is arranged to discuss feasibility and expectations.
  \item **Medication Adjustment:** Patients are instructed to discontinue blood-thinning medications (e.g., aspirin, NSAIDs, warfarin, novel oral anticoagulants) for a specified period before surgery, typically 5-7 days, to minimize bleeding risk.
  \item **NPO Status:** Patients must adhere strictly to NPO (nil per os) guidelines, typically refraining from solid food for at least 6-8 hours and clear liquids for 2 hours prior to surgery, to prevent aspiration during anesthesia.
  \item **Antibiotic Prophylaxis:** Prophylactic intravenous antibiotics (e.g., cefazolin) are administered within 60 minutes prior to skin incision to reduce the risk of surgical site infection.
  \item **Informed Consent:** A detailed discussion of the procedure, potential risks, benefits, alternatives, and expected outcomes is conducted, culminating in the patient providing fully informed consent.
\end{itemize}

Mastectomy, while a common and effective procedure, has specific contraindications and requires careful precautions in certain patient populations.

\textbf{Factors requiring optimization or a different treatment strategy:}
\begin{itemize}
  \item Uncontrolled systemic infection or sepsis generally requires stabilization before elective surgery, while urgent cancer or infection source control may still be necessary.
  \item Severe, uncorrectable coagulopathy or bleeding disorder that cannot be managed preoperatively, posing an unacceptable risk of hemorrhage.
  \item Moribund patient status (ASA V) where the risks of general anesthesia and major surgery far outweigh any potential benefits, and palliative care is more appropriate.
  \item Patient refusal after thorough counseling and understanding of the implications, as patient autonomy is paramount.
\item Distant metastatic disease with poor performance status where local treatment would not improve survival or quality of life, and systemic therapy is the primary focus.
\end{itemize}

**Relative Contraindications \& Precautions:**
\begin{itemize}
  \item Significant comorbidities (e.g., severe cardiac, pulmonary, or renal disease) that increase anesthetic and surgical risk; these require careful optimization and medical clearance prior to surgery.
  \item Active infection in the breast or surrounding tissues (e.g., cellulitis, abscess), which should be treated and resolved with antibiotics or drainage before elective surgery.
  \item Prior extensive radiation to the chest wall, which can compromise wound healing, increase the risk of skin flap necrosis, and complicate reconstructive options.
  \item Unrealistic patient expectations regarding cosmetic outcome, sensation, or recovery; thorough preoperative counseling and psychological support are crucial.
  \item Obesity (BMI $>$30), smoking, and uncontrolled diabetes, which are significant risk factors for wound complications (infection, dehiscence, flap necrosis) and should be optimized preoperatively.
  \item Psychological distress or severe body image concerns; patients should receive counseling and support from a multidisciplinary team to prepare for the physical and emotional changes.
  \item Pregnancy, especially in the first trimester, where surgical risks to the fetus are higher; surgery may be deferred or performed in the second trimester if necessary.
\end{itemize}

\section*{Surgical Technique \& Procedure}
The mastectomy procedure is performed under general anesthesia, with the patient typically positioned supine with the ipsilateral arm abducted on an arm board to allow optimal access to the axilla. The surgical field, extending from the clavicle to the umbilicus and from the sternum to the mid-axillary line, is prepped and draped in a sterile fashion.

The specific technique varies based on the type of mastectomy (e.g., simple, modified radical, skin-sparing, nipple-sparing), but the general steps include:

\begin{itemize}
  \item **1. Incision and Flap Elevation:** An elliptical incision is typically made to encompass the nipple-areolar complex and any prior biopsy scars, ensuring adequate skin removal for oncologic safety. For skin-sparing or nipple-sparing mastectomies, a smaller periareolar or inframammary incision is used to preserve the skin envelope or nipple-areolar complex. Skin flaps are then carefully raised using electrocautery or a scalpel, dissecting in the subcutaneous plane to a consistent thickness (typically 3-5 mm) to preserve vascularity while ensuring complete removal of glandular tissue.
  \item **2. Breast Tissue Dissection:** The dissection proceeds systematically, detaching the breast tissue from the pectoralis major fascia. The breast is freed from its superior (clavicle), inferior (inframammary fold), medial (sternum), and lateral (latissimus dorsi) attachments until the entire glandular tissue is mobilized. Care is taken to identify and ligate perforating vessels.
  \item \textbf{Axillary staging (if indicated):} Sentinel lymph node biopsy uses a tracer, dye, or both. The decision to perform axillary dissection depends on clinical stage, biopsy and sentinel-node findings, neoadjuvant treatment, and current breast-cancer guidance; it is not determined by frozen section in every case.
  \item **4. Hemostasis and Irrigation:** Thorough hemostasis is achieved using electrocautery or ligatures to minimize postoperative bleeding and hematoma formation. The surgical cavity is then irrigated with saline.
  \item **5. Drain Placement:** One or more closed-suction surgical drains (e.g., Jackson-Pratt drains) are placed in the breast cavity and/or axilla to prevent seroma formation by evacuating fluid. The drains are secured to the skin with sutures.
  \item **6. Closure:** The skin flaps are reapproximated and closed in layers. Absorbable sutures are used for the deep dermal layer to reduce tension, and non-absorbable sutures, staples, or surgical glue are used for the skin, often with a subcuticular closure for improved cosmetic outcome.
  \item **7. Reconstruction (if immediate):** If immediate breast reconstruction is planned, a plastic surgeon will perform the reconstructive phase following the oncologic mastectomy, utilizing either implant-based techniques or autologous tissue transfer (e.g., DIEP flap).
\end{itemize}

A sterile dressing is applied to the incision site.

\section*{Complications \& Risks}
Mastectomy, like any major surgical procedure, carries inherent risks, which can be broadly categorized into general surgical complications and procedure-specific risks.

\begin{itemize}
  \item **General Surgical Risks:**
    \begin{itemize}
      \item **Bleeding:** Intraoperative or postoperative hemorrhage, potentially requiring blood transfusion or reoperation for hematoma evacuation.
      \item **Infection:** Surgical site infection (SSI), ranging from superficial cellulitis to deep abscess formation, requiring antibiotics or drainage.
      \item **Anesthesia Complications:** Adverse reactions to anesthetic agents, respiratory or cardiac complications (e.g., myocardial infarction, arrhythmia), stroke, or, rarely, death.
      \item **Deep Vein Thrombosis (DVT) / Pulmonary Embolism (PE):** Formation of blood clots in the legs that can travel to the lungs, a potentially life-threatening complication, mitigated by prophylactic measures.
    \end{itemize}
  \item **Procedure-Specific Risks:**
    \begin{itemize}
      \item **Seroma:** Accumulation of sterile fluid under the skin flaps, often requiring aspiration or prolonged drain use. This is one of the most common complications, occurring in up to 30\% of patients.
      \item **Hematoma:** Collection of blood under the skin flaps, which may necessitate reoperation for evacuation to prevent infection or flap compromise.
      \item **Wound Dehiscence:** Partial or complete separation of the surgical incision, particularly in patients with poor wound healing, infection, or excessive tension.
      \item **Skin Flap Necrosis:** Death of skin tissue due to compromised blood supply, more common in smokers, diabetics, or after extensive skin elevation, potentially requiring debridement or further surgery.
      \item **Lymphedema:** Chronic swelling of the arm, hand, or chest wall on the operated side, particularly after axillary lymph node dissection, due to impaired lymphatic drainage. Incidence ranges from 5-30\%.
      \item **Nerve Injury:** Damage to specific nerves:
        \begin{itemize}
          \item **Intercostobrachial nerves:** Leading to numbness, tingling, or pain (neuroma) in the upper arm and axilla.
          \item **Long thoracic nerve:** Injury results in "winged scapula" due to paralysis of the serratus anterior muscle.
          \item **Thoracodorsal nerve:** Injury causes weakness of the latissimus dorsi muscle.
          \item **Medial/Lateral pectoral nerves:** Injury can lead to atrophy of the pectoralis muscles.
        \end{itemize}
      \item **Chronic Pain and Numbness:** Persistent pain (post-mastectomy pain syndrome) or altered sensation in the chest wall, axilla, or arm, which can significantly impact quality of life.
      \item **Cosmetic Dissatisfaction:** Concerns regarding breast asymmetry, scarring, or the overall aesthetic outcome, especially if reconstruction is not performed or is unsuccessful.
      \item **Phantom Breast Sensation:** A feeling that the breast is still present, or experiencing pain, itching, or tingling in the absent breast.
      \item **Implant-related complications (if reconstruction performed):** Infection, capsular contracture, rupture, malposition, or need for further surgery.
    \end{itemize}
\end{itemize}

\section*{Postoperative Care \& Recovery}
Immediately following a mastectomy, the patient is transferred to the Post-Anesthesia Care Unit (PACU) for close monitoring of vital signs, pain control, and assessment for immediate complications such as bleeding or respiratory distress. Once stable, they are moved to a surgical ward. Pain management is a priority, typically involving a multimodal approach with oral or intravenous analgesics. Surgical drains, usually one or two, are present to collect fluid and prevent seroma formation; patients and their caregivers are taught how to manage and record drain output. Early ambulation is strongly encouraged to prevent complications like deep vein thrombosis and pulmonary embolism. The typical hospital stay ranges from one to three days, depending on the extent of surgery, whether immediate reconstruction was performed, and the patient's overall recovery trajectory.

During the first 1-2 weeks post-discharge, patients focus on wound care, drain management, and gradual resumption of light activities. Drains are usually removed when output consistently falls below a certain threshold (e.g., 20-30 mL per 24 hours) for several consecutive days, typically within 1-3 weeks. Patients are advised to avoid heavy lifting (generally anything over 5-10 pounds), strenuous activities, and repetitive arm movements on the operated side to prevent wound complications, promote healing, and minimize the risk of seroma. Physical therapy may be initiated early to restore full range of motion in the arm and shoulder, especially after axillary dissection, and to prevent "cording" (axillary web syndrome). Follow-up appointments with the surgeon are scheduled to monitor wound healing, remove drains, and discuss the definitive pathology results. Long-term recovery involves managing potential lymphedema, addressing body image concerns, and proceeding with any necessary adjuvant therapies or delayed reconstructive procedures. Full recovery of strength and sensation can take several months to a year.

During the mastectomy procedure, the surgeon meticulously documents the gross findings, including the size and location of the tumor, its relationship to the skin and pectoralis fascia, and the presence of any suspicious axillary lymph nodes. The resected breast tissue, along with any removed lymph nodes, is then sent to pathology for detailed microscopic analysis. The pathology report is the definitive document for cancer staging and guiding subsequent adjuvant treatments.

Key information in the pathology report includes the precise tumor type (e.g., invasive ductal carcinoma, invasive lobular carcinoma, ductal carcinoma in situ), its exact size (measured in millimeters), and histological grade (e.g., Nottingham grade 1, 2, or 3), which indicates the aggressiveness of the cancer cells. The report will also specify the presence or absence of lymphovascular invasion (cancer cells in small blood or lymphatic vessels) and perineural invasion. Crucially, the status of surgical margins is reported, indicating whether the tumor was completely removed with a rim of healthy tissue (negative margins) or if cancer cells extend to the edge of the resected specimen (positive margins), which may necessitate further treatment. For invasive cancers, the pathology report details the status of hormone receptors (estrogen receptor [ER] and progesterone receptor [PR]) and HER2 (human epidermal growth factor receptor 2) expression, which are critical biomarkers for determining eligibility for endocrine therapy and HER2-targeted therapies, respectively. The number of lymph nodes removed and the number of nodes positive for metastatic disease are also reported, which is a major factor in cancer staging (e.g., pN0, pN1) and prognosis. For prophylactic mastectomies, the pathology report confirms the absence of malignancy or may identify incidental findings such as atypical hyperplasia or lobular carcinoma in situ, which are important for future risk assessment.

A normal and successful postoperative course following a mastectomy is characterized by several key indicators of healing and recovery. Patients can expect initial pain, which should be well-controlled with prescribed analgesics and gradually diminish over the first few weeks. The surgical incision should heal cleanly, without signs of infection (redness, warmth, pus), significant dehiscence, or extensive skin flap necrosis. Drains will typically remain in place for 1-3 weeks, with output gradually decreasing until they can be safely removed. Seroma formation, while common, should be manageable, often resolving with drain removal or occasional aspiration.

Patients should progressively regain full or near-full range of motion in the ipsilateral arm and shoulder, especially with adherence to physical therapy exercises, and without significant lymphedema. While some numbness or altered sensation in the chest wall and arm is common and often permanent due to nerve transection, severe or debilitating chronic pain should be rare. For therapeutic mastectomies, a successful course also implies that the pathology report confirms clear surgical margins and provides comprehensive information to guide appropriate adjuvant therapies. Ultimately, a normal course leads to a patient who is recovering well physically and emotionally, able to resume most daily activities, and prepared for any necessary subsequent treatments or reconstructive procedures, with a good overall quality of life.

Postoperative care and long-term follow-up after mastectomy are crucial for monitoring recovery, detecting potential recurrence, and managing any chronic complications.

\begin{itemize}
  \item **Surgical Follow-up:** Initial visits with the surgical oncologist typically occur within 1-3 weeks post-surgery for drain removal, wound check, and detailed discussion of the pathology results. Subsequent visits monitor healing and address any immediate concerns.
  \item **Oncologic Surveillance:** Regular follow-up with the medical oncologist and/or radiation oncologist is crucial. This includes clinical examinations of the chest wall, axilla, and contralateral breast every 3-6 months for the first 2-3 years, then annually for at least 5-10 years.
  \item **Imaging:** Annual mammography of the contralateral breast is recommended for all patients. Imaging of the chest wall or reconstructed breast (e.g., ultrasound, MRI) may be performed if clinically indicated by new symptoms or suspicious findings.
  \item **Blood Tests:** Routine tumor markers or blood work are generally not recommended for surveillance in asymptomatic patients but may be ordered if there are specific concerns for recurrence.
  \item **Rehabilitation/Physical Therapy:** Many patients benefit from ongoing physical therapy to regain full range of motion, strength, and address issues like cording (axillary web syndrome) or lymphedema. Referral to a certified lymphedema therapist is important if swelling develops.
  \item **Psychosocial Support:** Ongoing support for body image, emotional well-being, and sexual health is often vital, with referrals to support groups, psychologists, or specialized counselors as needed.
  \item **Warning Signs:** Patients are educated on warning signs of local recurrence (e.g., new lumps, skin changes, persistent pain, swelling in the chest wall or axilla) and symptoms of distant metastasis (e.g., new bone pain, persistent cough, headaches, unexplained weight loss) to report promptly to their healthcare team.
\end{itemize}
Long-term follow-up focuses on overall health, screening for new primary cancers (especially in the contralateral breast for those who did not undergo bilateral mastectomy), and managing any chronic side effects of treatment.

\section*{Outcomes \& Prognosis}
Mastectomy remains a common surgical procedure globally, primarily driven by the high incidence of breast cancer. Breast cancer is the most frequently diagnosed cancer among women worldwide, accounting for approximately 1 in 8 cancer diagnoses. In the United States, approximately 1 in 8 women (13\%) will develop invasive breast cancer over their lifetime. While breast-conserving surgery (BCS) is increasingly utilized and often preferred for early-stage disease, mastectomy rates vary significantly by region, patient demographics, tumor characteristics, and access to reconstructive surgery. Approximately 30-40\% of women undergoing surgery for early-stage breast cancer in the U.S. opt for mastectomy, with rates potentially higher in certain populations (e.g., younger women, those with genetic mutations) or for specific indications like multifocal disease, large tumor size, or inflammatory breast cancer. Prophylactic mastectomy, while less common, is performed in a growing number of high-risk individuals, particularly those with BRCA1/2 mutations, where it can reduce the lifetime risk of breast cancer by over 90\%. Morbidity associated with mastectomy, such as seroma and lymphedema, affects a substantial proportion of patients, with lymphedema incidence ranging from 5-30\% depending on the extent of axillary surgery. Mortality directly attributable to mastectomy is very low, typically less than 0.1-0.5\%, primarily related to anesthesia or major surgical complications in medically fragile patients.

The outcomes and prognosis following mastectomy are highly dependent on the underlying indication, particularly the stage and biological characteristics of the breast cancer. For early-stage invasive breast cancer, mastectomy offers excellent local control, with recurrence rates in the ipsilateral chest wall typically less than 5-10\% over 10 years. Survival outcomes are generally very favorable, with 5-year survival rates exceeding 90\% for stage I and II disease, often comparable to breast-conserving surgery with radiation, as demonstrated by landmark trials such as NSABP B-04 and B-06. Prognostic factors include tumor size, nodal status (the most significant factor), histological grade, hormone receptor status (ER/PR), and HER2 expression. Patients with negative lymph nodes and favorable tumor biology have the best prognosis.

For prophylactic mastectomy in high-risk individuals, the procedure is highly effective, reducing the lifetime risk of breast cancer by 90-95\%. While it significantly lowers the risk, it does not eliminate it entirely due to the potential for microscopic residual breast tissue. Functional outcomes, such as arm range of motion, are generally good, especially with early physical therapy and adherence to rehabilitation protocols. However, long-term issues like lymphedema, chronic pain, and altered sensation can impact quality of life for a subset of patients. Psychological adjustment and body image satisfaction vary widely, with reconstructive surgery playing a significant role in improving patient satisfaction and overall quality of life. Overall, mastectomy remains a highly effective and often life-saving procedure, with continuous advancements in surgical technique, adjuvant therapies, and reconstructive options further improving patient outcomes and prognosis.

\section*{References}
\begin{enumerate}
  \item MedlinePlus. Mastectomy. U.S. National Library of Medicine; 2024. Available from: \url{https://medlineplus.gov/mastectomy.html}
  \item Bland KI, Copeland EM, Klimberg VS, Gradishar WJ. The Breast: Comprehensive Management of Benign and Malignant Diseases. 5th ed. Elsevier; 2018.
  \item National Comprehensive Cancer Network (NCCN). NCCN Clinical Practice Guidelines in Oncology: Breast Cancer. Version 1.2024. Available from: \url{https://www.nccn.org/guidelines/category_1}
  \item American Cancer Society. Mastectomy. 2024. Available from: \url{https://www.cancer.org/cancer/types/breast-cancer/treatment/surgery-for-breast-cancer/mastectomy.html}
\end{enumerate}

\clearpage
